\documentclass[9pt]{beamer}
\usetheme{Madrid}
\usepackage[utf8]{inputenc}
\usepackage[english]{babel}
\usepackage{amsmath}
\usepackage{amsfonts}
\usepackage{amssymb}
\usepackage{graphicx}
\usepackage{hyperref}
\usepackage{adjustbox}
\usepackage{pgfpages}

\usepackage{array}  % image alignment in tabular
\author{Pietro Mascolo}
\title{Sequence modeling for product classification}
\logo{\includegraphics[height=0.45cm]{flow_resources/flow-logo-icon-&-type-color.png}}
%\setbeamercovered{transparent} 
%\setbeamertemplate{navigation symbols}{} 
%\logo{} 
\institute{\href{http://www.flow.io}{Flow Comerce}} 
%\date{} 
%\subject{} 

\setbeamertemplate{note page}[plain]
\setbeameroption{show notes on second screen=right}

% Flow palette
\definecolor{FlowBlue}{RGB}{00,81,178}
\definecolor{VibrantBlue}{RGB}{71,145,255}
\definecolor{Steel}{RGB}{128,128,128}
\definecolor{Silver}{RGB}{191,191,191}
\definecolor{LightSilver}{RGB}{220,220,220}
\definecolor{VeryLightSilver}{RGB}{240,240,240}
\definecolor{Accent1}{RGB}{85,218,255}
\definecolor{Accent2}{RGB}{249,119,98}
\definecolor{Accent3}{RGB}{250,221,5}
\definecolor{Accent4}{RGB}{82,212,253}
\definecolor{AccentAlert}{RGB}{249,50,50}

% More info on sections and parts of beamer here: http://www.cpt.univ-mrs.fr/~masson/latex/Beamer-appearance-cheat-sheet.pdf

% general
\setbeamercolor{palette primary}{bg=FlowBlue}
\setbeamercolor{palette secondary}{bg=VibrantBlue}
\setbeamercolor{palette tertiary}{bg=Steel}
\setbeamercolor{palette quaternary}{bg=Accent1}
\setbeamercolor{palette sidebar primary}{bg=Accent1}
\setbeamercolor{palette sidebar secondary}{bg=Accent2}
\setbeamercolor{palette sidebar tertiary}{bg=Accent2}
\setbeamercolor{palette sidebar quaternary}{bg=Accent4}

% hyperlinks
\setbeamercolor{button}{use=Accent4}
\setbeamercolor{button border}{use=Accent4}

% structure
\setbeamercolor{structure}{fg=FlowBlue}
\setbeamercolor{local structure}{fg=Accent4}
\setbeamercolor{alerted text}{fg=AccentAlert, bg=LightSilver}

% block
\setbeamercolor{block title}{bg=FlowBlue}
\setbeamercolor{block body}{bg=VeryLightSilver}

% titles
\setbeamercolor{title}{bg=FlowBlue, fg=white}
\setbeamercolor{frametitle}{bg=FlowBlue, fg=white}

%itemization
%\setbeamercolor{itemize item}{bg=FlowBlue}
%\setbeamercolor{itemize subitem}{bg=FlowBlue}
%\setbeamercolor{itemize subsubitem}{bg=FlowBlue}
%\setbeamercolor{enumerate item}{bg=FlowBlue}
%\setbeamercolor{enumerate subitem}{bg=FlowBlue}
%\setbeamercolor{enumerate subsubitem}{bg=FlowBlue}
%\setbeamertemplate{itemize item}{\color{FlowBlue}}
%\setbeamertemplate{enumerate item}{\color{FlowBlue}}

%\setbeamerfont{section number projected}{family=\rmfamily,series=\bfseries,size=\normalsize}
\setbeamercolor{section number projected}{bg=FlowBlue,fg=white}
\setbeamercolor{subsection number projected}{bg=FlowBlue}

\setbeamercolor{structure}{bg=FlowBlue}

% margins
\setbeamersize{text margin left=10mm,text margin right=10mm}
\setbeamertemplate{caption}{\raggedright\insertcaption\par}


\begin{document}

%%%%%%%%%%%%%%%%%%%%%%%%%%%%%%%%%%%%%%%%%%%%
\begin{frame}
	\titlepage
\end{frame}


%%%%%%%%%%%%%%%%%%%%%%%%%%%%%%%%%%%%%%%%%%%%
\begin{frame}{This presentation}
	\tableofcontents
\end{frame}


\section{Problem statement}
%%%%%%%%%%%%%%%%%%%%%%%%%%%%%%%%%%%%%%%%%%%%
\begin{frame}{Introduction}
	\begin{block}{Problem statement}
		We want to classify a generic product into an HS6 code and a brief description in order to determine duty rates for a given origin and destination.
	\end{block}
	\vspace{0.5cm}
	
	\pause
	\begin{center}
		\begin{tabular}{ c  c m{5cm} }
			\begin{minipage}{4cm}
			\begin{figure}
					\includegraphics[width=4cm]{imgs/boot.png}
					 \tiny{A classically styled pair of leather Chelsea boots, featuring elasticated side panels, heel pull tabs and finished with a subtle...}
				\end{figure}
			\end{minipage} \pause
			& $\longrightarrow$ 
			& $$\left\{ 
				\begin{array}{l} 
					\text{description: womens leather boot} \\
					\\
					\text{HS6: 640391}
				\end{array} 
			\right.$$ \\
		\end{tabular}
	\end{center}
\end{frame}


\section{First approach}
%%%%%%%%%%%%%%%%%%%%%%%%%%%%%%%%%%%%%%%%%%%%
\begin{frame}{Our initial approach}
	\begin{center}
		\includegraphics[width=0.4\textwidth,height=!,trim={0 0 30cm 5cm},clip]{imgs/houses.jpg}
	\end{center}
\end{frame}

%%%%%%%%%%%%%%%%%%%%%%%%%%%%%%%%%%%%%%%%%%%%
\begin{frame}{Initial models}
	\begin{itemize}
		\item TFIDF encoded representation of text;
		\item one model per organization (20-30);
		\item reasonably accurate.
	\end{itemize}
	\pause
	\vspace{1cm}
	\begin{center}
		\Large{\textbf{BUT...}}
	\end{center}
\end{frame}

%%%%%%%%%%%%%%%%%%%%%%%%%%%%%%%%%%%%%%%%%%%%
\begin{frame}{Little to no generalization}
	\begin{center}
		\begin{figure}
			\includegraphics[width=0.6\textwidth,height=!]{imgs/huffedandpuffed.png}
			\caption{\textit{...and he huffed and he puffed and he blew the house in...}}
		\end{figure}
	\end{center}
\end{frame}

%%%%%%%%%%%%%%%%%%%%%%%%%%%%%%%%%%%%%%%%%%%%
\begin{frame}{Problems!!}
	\begin{itemize}
		\item TFIDF encoding doesn't capture sequence relationship well! %\note{TFIDF can represent some n-gram features, but in a very limited way.}
		\item Many models to keep track of!
		\item Hard performance evaluation!
		\item Some organizations have very little data to work on!
		\item Some organizations have a very limited catalogue (1 or 2 types of items)!
	\end{itemize}

	\pause
	\begin{center}
		\includegraphics[width=0.2\textwidth]{imgs/sad.png}
	\end{center}
\end{frame}


\section{A new approach}
%%%%%%%%%%%%%%%%%%%%%%%%%%%%%%%%%%%%%%%%%%%%
\begin{frame}{A new approach}
	\begin{center}
		\includegraphics[width=0.4\textwidth,height=!,trim={16cm 0 17.5cm 5cm},clip]{imgs/houses.jpg}
	\end{center}
\end{frame}


%%%%%%%%%%%%%%%%%%%%%%%%%%%%%%%%%%%%%%%%%%%%
\begin{frame}{Sequence modeling}
	We want to:
	\begin{itemize}
		\item Have a single general model;
		\item capture information regarding patterns in sequences of words;
		\item be able add a new organization seamlessly (if we have seen the kind of products they sell).
	\end{itemize}
	\vspace{0.5cm}
	These requirements suggest that we need a more sophisticated method than what we are using.
	
	\pause
	\vspace{1cm}
	\begin{center}
		\Large{\textbf{Enter Neural Networks...}}
	\end{center}
	
\end{frame}

\subsection{Neural networks}
%%%%%%%%%%%%%%%%%%%%%%%%%%%%%%%%%%%%%%%%%%%%
\begin{frame}{A primer on neural networks}

\end{frame}

%%%%%%%%%%%%%%%%%%%%%%%%%%%%%%%%%%%%%%%%%%%%
\begin{frame}{How Neural Networks learn}

\end{frame}


\subsection{Recurrent networks and sequence modeling}

%%%%%%%%%%%%%%%%%%%%%%%%%%%%%%%%%%%%%%%%%%%%
\begin{frame}{Sequence learning - Recurrent Neural Networks}

\end{frame}

%%%%%%%%%%%%%%%%%%%%%%%%%%%%%%%%%%%%%%%%%%%%
\begin{frame}{Short Term memory problem - vanishing gradient}

\end{frame}

%%%%%%%%%%%%%%%%%%%%%%%%%%%%%%%%%%%%%%%%%%%%
\begin{frame}{LSTM}

\end{frame}

%%%%%%%%%%%%%%%%%%%%%%%%%%%%%%%%%%%%%%%%%%%%
\begin{frame}{GRU}

\end{frame}

\section{Our classification model}

%%%%%%%%%%%%%%%%%%%%%%%%%%%%%%%%%%%%%%%%%%%%
\begin{frame}{The model}

\end{frame}

\subsection{Data preparation}
%%%%%%%%%%%%%%%%%%%%%%%%%%%%%%%%%%%%%%%%%%%%
\begin{frame}{The data}

\end{frame}


\subsection{Evaluation}
%%%%%%%%%%%%%%%%%%%%%%%%%%%%%%%%%%%%%%%%%%%%
\begin{frame}{Results and evaluation}

\end{frame}

%%%%%%%%%%%%%%%%%%%%%%%%%%%%%%%%%%%%%%%%%%%%
\begin{frame}{Problems??}
	\begin{center}
		\begin{figure}
			\includegraphics[width=0.6\textwidth,height=!]{imgs/huffedandpuffed.png}
			\caption{\textit{...and he huffed and he puffed, \textbf{but he hasn't blown the house in... yet...}}}
		\end{figure}
	\end{center}
\end{frame}

\section{Demo}
%%%%%%%%%%%%%%%%%%%%%%%%%%%%%%%%%%%%%%%%%%%%
\begin{frame}{Let's see it in action}

\begin{center}
	\Huge{\href{http://localhost:8888}{Demo!}}
\end{center}

\end{frame}




%%%%%%%%%%%%%%%%%%%%%%%%%%%%%%%%%%%%%%%%%%%%
\begin{frame}{Conclusions and future work}
	\begin{center}
		\includegraphics[width=0.4\textwidth,height=!,trim={29cm 0 0 5.2cm},clip]{imgs/houses.jpg}
	\end{center}
\end{frame}
\end{document}